% ===========================================================================
\section{Formalising ``lossless''}\label{sec:lossless}
% ===========================================================================

\subsection{Divergences, distillation risk and task-level excess risk}

\begin{definition}[$f$-divergences used throughout]\label{def:fdiv}
For convex $f:(0,\infty)\to\R$ with $f(1)=0$, set
$D_f(P\|Q)=\sum_{y}Q(y)\,f\!\big(P(y)/Q(y)\big)$. We use the conventions $0f(0/0)=0$ and
$Qf(P/0)=P\,f^*(0)$, where $f^*(0):=\lim_{u\to\infty}f(u)/u$. The instances we use are:
\begin{center}\small
\begin{tabular}{@{}llll@{}}
\toprule
divergence & $f(u)$ & $f''(1)$ & bounded?\\\midrule
forward KL, $\KL(P\|Q)$ & $u\log u$ & $1$ & no\\
reverse KL, $\KL(Q\|P)$ & $-\log u$ & $1$ & no\\
total variation $\TV=\frac12\|P-Q\|_1$ & $\frac12|u-1|$ & -- & $\le1$\\
Hellinger$^2$, $\Hel^2=\frac12\sum(\sqrt P-\sqrt Q)^2$ & $\frac12(\sqrt u-1)^2$ & $\frac14$ & $\le 1$\\
$\chi^2(P\|Q)$ & $(u-1)^2$ & $2$ & no\\
$\JSD_\beta(P\|Q)$ & $\beta u\log u-m\log m$, $m=\beta u+1-\beta$ & $\beta(1-\beta)$ & $\le h(\beta)$\\\bottomrule
\end{tabular}
\end{center}
Here $\JSD_\beta(P\|Q):=\beta\,\KL(P\|M_\beta)+(1-\beta)\,\KL(Q\|M_\beta)$ with $M_\beta=\beta P+(1-\beta)Q$.
This follows the convention of \citet{agarwal2024gkd}; $\beta=\frac12$ is the classical
Jensen--Shannon divergence \citep{lin1991divergence}. Conventions differ between papers (some swap
$\beta$ and $1-\beta$), so we fix this one. A direct computation gives
$f_\beta''(u)=\beta(1-\beta)/(u\,m)>0$. The Rényi divergence of order $\alpha>1$ is
$D_\alpha(P\|Q)=\frac{1}{\alpha-1}\log\sum_y P(y)^\alpha Q(y)^{1-\alpha}$ \citep{vanerven2014renyi}.
\end{definition}

\begin{definition}[Distillation risk; $(\eps,\mathsf D,\mu)$-lossless]\label{def:risk}
For a divergence $\mathsf D$ on $\Delta(\cY)$ and an input law $\mu$, the \emph{sequence-level
distillation risk} is
$\cL_{\mathsf D}(\theta;\mu):=\E_{x\sim\mu}\,\mathsf D\big(\Pseq\,\|\,\Qseq\big)$.
The student is \emph{$(\eps,\mathsf D,\mu)$-lossless} if $\cL_{\mathsf D}(\theta;\mu)\le\eps$.
For a family of prefix laws $\nu=(\nu_t)_{t\le T}$, the \emph{token-level risk} is
$\ell_{\mathsf D}(\theta;\nu):=\sum_{t=1}^T\E_{s\sim\nu_t}\mathsf D\big(p_t(\cdot\mid s)\|q_t(\cdot\mid s)\big)$.
\end{definition}

\begin{definition}[Task-level excess risk]\label{def:task}
A task is a measurable reward $r:\cX\times\cY\to[0,1]$; a loss $\ell$ is handled as $r=1-\ell$.
Its excess risk under sampling is
$\Delta_r(\theta;\mu):=\E_{x\sim\mu}\big[\E_{y\sim\Pseq}r(x,y)-\E_{y\sim\Qseq}r(x,y)\big]$.
For a class $\cR$ of tasks, $\Delta_{\cR}(\theta;\mu):=\sup_{r\in\cR}|\Delta_r(\theta;\mu)|$.
\end{definition}

\subsection{Relations between the definitions}

\begin{proposition}[TV is the worst-case task regret]\label{prop:tv-task}
\Pv
\begin{enumerate}[label=(\alph*),nosep]
\item $\sup_{r:\cX\times\cY\to[0,1]}\Delta_r(\theta;\mu)=\sup_{r}|\Delta_r(\theta;\mu)|=\E_{x\sim\mu}\TV(\Pseq,\Qseq)$.
\item If the task sees $y$ only through a statistic $\phi(x,y)$, for example the final answer, so that
$r=g(x,\phi(x,y))$, then
$\sup_g|\Delta_r|=\E_{x\sim\mu}\TV(\phi_\#\Pseq,\phi_\#\Qseq)\le\E_{x\sim\mu}\TV(\Pseq,\Qseq)$.
\item For each $x$, the minimum error probability of any test that tells, from one draw, whether the
draw came from $\Pseq$ or from $\Qseq$ (equal priors) is $\frac12\big(1-\TV(\Pseq,\Qseq)\big)$.
\end{enumerate}
\end{proposition}
\begin{proof}
(a) For fixed $x$, $\sum_y r(x,y)(\Pseq(y)-\Qseq(y))\le\sum_y(\Pseq(y)-\Qseq(y))_+=\TV$, with equality at
$r=\ind\{\Pseq>\Qseq\}$. This pointwise maximiser is measurable because $\cX\times\cY$ is countable.
Taking $1-r$ covers the absolute value. Integrate over $x\sim\mu$.
(b) Apply (a) to the push-forward laws, then use data processing for $\TV$.
(c) This is Le Cam's two-point identity: the Bayes test is $\ind\{\Pseq>\Qseq\}$.
\end{proof}

\begin{proposition}[Pinsker, Bretagnolle--Huber, Hellinger]\label{prop:pinsker}
\Kn For any $P,Q$ and $\mathsf K\in\{\KL(P\|Q),\KL(Q\|P)\}$:
(i)~$\TV\le\sqrt{\mathsf K/2}$ \citep{pinsker1964information,tsybakov2009introduction};
(ii)~$\TV\le\sqrt{1-e^{-\mathsf K}}$ \citep{bretagnolle1979estimation}, which is the informative bound when
$\mathsf K>2$;
(iii)~$\Hel^2\le\TV\le\Hel\sqrt{2-\Hel^2}$;
(iv)~$\mathsf K\ge-2\log(1-\Hel^2)\ge2\Hel^2$.
\end{proposition}
\begin{proof}[Proof of (iv)]
By Jensen, $\KL(P\|Q)=-2\,\E_P\log\sqrt{Q/P}\ge-2\log\E_P\sqrt{Q/P}=-2\log(1-\Hel^2)$. The same argument
applies with $P$ and $Q$ exchanged. The other parts are textbook; see \citet[Lemmas~2.3--2.6]{tsybakov2009introduction}.
\end{proof}

\begin{lemma}[$f$-divergences are at most linear in TV]\label{lem:vajda}
\Pv For convex $f$ with $f(1)=0$, $D_f(P\|Q)\le\big(f(0)+f^*(0)\big)\,\TV(P,Q)$.
\end{lemma}
\begin{proof}
Write $u=P/Q$. On $[0,1]$, convexity gives $f(u)\le(1-u)f(0)+uf(1)=(1-u)f(0)$. For $u>1$ the slope
$(f(u)-f(1))/(u-1)$ is non-decreasing in $u$ and tends to $f^*(0)$, so $f(u)\le(u-1)f^*(0)$. Hence
$D_f\le f(0)\sum_y(Q-P)_++f^*(0)\sum_y(P-Q)_+=(f(0)+f^*(0))\TV$. This is classical; see
\citet{sason2016fdivergence}.
\end{proof}

\begin{proposition}[Properties of $\JSD_\beta$]\label{prop:jsd}
\Pv\Ck For all $P,Q$ and $\beta\in(0,1)$:
\begin{enumerate}[label=(\roman*),nosep]
\item $2\beta(1-\beta)\,\TV^2\le\JSD_\beta(P\|Q)\le h(\beta)\,\TV\le h(\beta)$;
\item $\JSD_\beta(P\|Q)\le\beta(1-\beta)\big[\KL(P\|Q)+\KL(Q\|P)\big]$;
\item if $\operatorname{supp}P=\operatorname{supp}Q$ is finite, then
$\JSD_\beta/\beta\to\KL(P\|Q)$ as $\beta\to0$ and $\JSD_\beta/(1-\beta)\to\KL(Q\|P)$ as $\beta\to1$.
\end{enumerate}
\end{proposition}
\begin{proof}
(i) Since $\TV(P,M_\beta)=(1-\beta)\TV(P,Q)$ and $\TV(Q,M_\beta)=\beta\,\TV(P,Q)$, Pinsker applied to each
term gives $\JSD_\beta\ge2\beta(1-\beta)^2\TV^2+2(1-\beta)\beta^2\TV^2$. For the upper bound, apply
\cref{lem:vajda} with $f_\beta(0)=(1-\beta)\log\frac1{1-\beta}$ and
$f_\beta^*(0)=\lim_u[\beta\log u-\beta\log(\beta u)]=\beta\log\frac1\beta$, whose sum is $h(\beta)$.
(ii) KL is convex in its second argument, so $\KL(P\|\beta P+(1-\beta)Q)\le(1-\beta)\KL(P\|Q)$ and
$\KL(Q\|M_\beta)\le\beta\,\KL(Q\|P)$.
(iii) As $\beta\to0$, $M_\beta\to Q$, so $\KL(P\|M_\beta)\to\KL(P\|Q)$. Also
$\KL(Q\|Q+\beta(P-Q))=\frac{\beta^2}2\chi^2(P\|Q)+O(\beta^3)$ on a finite common support. Hence
$\JSD_\beta=\beta\KL(P\|Q)+O(\beta^2)$. The case $\beta\to1$ is symmetric.
\end{proof}

\begin{proposition}[No global comparison; local equivalence]\label{prop:local}
\Pv\Ck
\begin{enumerate}[label=(\roman*),nosep]
\item Forward and reverse KL are not comparable in general. For $P=\mathrm{Bern}(\delta)$ and
$Q=\mathrm{Bern}(0)$, $\KL(P\|Q)=\infty$ while $\KL(Q\|P)=\log\frac1{1-\delta}\to0$. Exchanging the
roles gives the opposite.
\item Suppose $a\le P(y)/Q(y)\le b$ for all $y$ (necessarily $a\le1\le b$). For every twice
differentiable $f$ with $f(1)=0$,
\[
\tfrac12\inf_{[a,b]}f''\cdot\chi^2(P\|Q)\;\le\;D_f(P\|Q)\;\le\;\tfrac12\sup_{[a,b]}f''\cdot\chi^2(P\|Q).
\]
In particular $\frac{\chi^2}{2b}\le\KL(P\|Q)\le\frac{\chi^2}{2a}$ and
$\frac{\chi^2}{2b^2}\le\KL(Q\|P)\le\frac{\chi^2}{2a^2}$. Hence $\KL(P\|Q)\le\frac{b^2}{a}\KL(Q\|P)$ and
$\KL(Q\|P)\le\frac{b}{a^2}\KL(P\|Q)$. Also $\JSD_\beta=\frac{\beta(1-\beta)}2\chi^2(1+O(b-a))$.
\end{enumerate}
\end{proposition}
\begin{proof}
(ii) Replacing $f$ by $\tilde f(u)=f(u)-f'(1)(u-1)$ leaves $D_f$ unchanged, because $\sum_yQ(u-1)=0$.
Taylor's theorem gives $\tilde f(u)=\frac12f''(\xi_u)(u-1)^2$ with $\xi_u\in[a,b]$. Multiply by $Q$ and sum.
For forward KL, $f''(u)=1/u$; for reverse KL, written as $D_{-\log}(P\|Q)$, $f''(u)=1/u^2$.
\end{proof}

\begin{remark}[What \cref{prop:local} means for ``near-lossless'']\label{rem:local}
In the near-lossless regime with bounded likelihood ratios, every smooth $f$-divergence equals
$\frac12f''(1)\chi^2$ to leading order. So forward KL, reverse KL and $\JSD_\beta$ (rescaled by
$\beta(1-\beta)$) \emph{agree} when the student is nearly lossless. The choice of divergence matters
only (a)~in the lossy, capacity-limited regime, (b)~through tails and support mismatch, where likelihood
ratios are unbounded, and (c)~through \emph{which prefix law} the per-token divergence is averaged over
(\cref{sec:chain}).
\end{remark}

\begin{proposition}[Coverage versus precision: event bounds]\label{prop:events}
\Pv\Ck For every event $A\subseteq\cY$:
\begin{enumerate}[label=(\roman*),nosep]
\item (coverage) $P(A)\le\dfrac{\KL(P\|Q)+\log2}{\log(1/Q(A))}$;
\item (precision) $Q(A)\le\dfrac{\KL(Q\|P)+\log2}{\log(1/P(A))}$;
\item (Rényi, $\alpha>1$) $Q(A)\le\big(e^{D_\alpha(Q\|P)}P(A)\big)^{(\alpha-1)/\alpha}$;
\item (additive) $|P(A)-Q(A)|\le\TV(P,Q)$.
\end{enumerate}
\end{proposition}
\begin{proof}
(ii) By data processing onto $\{A,A^c\}$, $\KL(Q\|P)\ge\mathrm{kl}(Q(A)\|P(A))$, and
$\mathrm{kl}(a\|b)\ge a\log\frac1b-h(a)\ge a\log\frac1b-\log2$. Part (i) is (ii) with the roles
exchanged. (iii) By Hölder with exponents $\alpha$ and $\alpha/(\alpha-1)$,
$Q(A)=\sum_AP\,(Q/P)\le P(A)^{(\alpha-1)/\alpha}\big(\sum_yP(Q/P)^\alpha\big)^{1/\alpha}$, and
$\sum_yP(Q/P)^\alpha=e^{(\alpha-1)D_\alpha(Q\|P)}$ \citep[Prop.~10]{mironov2017renyi}.
(iv) is the definition of TV.
\end{proof}
\noindent
Forward KL guarantees \emph{coverage}: the student cannot give negligible probability to events the
teacher finds likely. Reverse KL, and Rényi divergences of $Q$ from $P$, guarantee \emph{precision}: the
student cannot put appreciable mass on events the teacher considers essentially impossible (for example
hallucinated or unsafe continuations). TV controls neither multiplicatively. A student at
$\TV=10^{-2}$ may emit a teacher-impossible output with probability $10^{-2}$, which is invisible to TV
but catastrophic at deployment scale.

\begin{proposition}[Greedy decoding needs margins]\label{prop:margin}
\Pv Let $a^\star(s)=\argmax_a p(a\mid s)$ have margin
$\gamma(s)=p(a^\star\mid s)-\max_{a\ne a^\star}p(a\mid s)$. If $\TV(p_t(\cdot\mid s),q_t(\cdot\mid s))<\gamma(s)/2$
at every state $s$ on the teacher's greedy path from $x$, then the student's greedy decoding of $x$
equals the teacher's. The constant is tight: moving mass $\gamma/2$ from $a^\star$ to the runner-up
creates a tie. \Cref{sec:decoding} turns this condition into sharp bounds on the probability that the
outputs differ.
\end{proposition}
\begin{proof}
$|p(a\mid s)-q(a\mid s)|\le\TV$ for every $a$, so
$q(a^\star\mid s)-q(a\mid s)\ge\gamma(s)-2\TV>0$. Induct along the path.
\end{proof}

\subsection{Sequence level versus token level}\label{sec:chain}

\begin{theorem}[Chain rules: each KL direction has its own prefix law]\label{thm:chain}
\Pv\Ck For every $x$,
\begin{align}
\KL(\Pseq\|\Qseq)&=\sum_{t=1}^T\E_{s_t\sim\Pseq}\,\KL\big(p_t(\cdot\mid s_t)\,\|\,q_t(\cdot\mid s_t)\big),\label{eq:chain-fwd}\\
\KL(\Qseq\|\Pseq)&=\sum_{t=1}^T\E_{s_t\sim\Qseq}\,\KL\big(q_t(\cdot\mid s_t)\,\|\,p_t(\cdot\mid s_t)\big).\label{eq:chain-rev}
\end{align}
After averaging over $x\sim\cD$: $\E_\cD\KL(P\|Q_\theta)=T\,\eoff^{\KL}$ (teacher prefixes) and
$\E_\cD\KL(Q_\theta\|P)=T\,\eon^{\mathrm{rKL}}$ (student prefixes).
\end{theorem}
\begin{proof}
$\log\frac{\Pseq(y)}{\Qseq(y)}=\sum_t\log\frac{p_t(y_t\mid s_t)}{q_t(y_t\mid s_t)}$. Take $\E_{y\sim\Pseq}$ and
condition each summand on $s_t$ (tower property). The reverse identity is the same computation with
$y\sim\Qseq$.
\end{proof}

\begin{remark}\label{rem:pairings}
Supervised distillation on teacher-generated sequences with soft labels at every position (the
token-level analogue of \citet{kim2016sequence}) minimises \eqref{eq:chain-fwd} \emph{exactly}. On-policy
reverse-KL distillation, as in the objectives of \citet{gu2024minillm} and \citet{agarwal2024gkd} with
$\lambda=1$ and reverse KL, minimises \eqref{eq:chain-rev} exactly, provided the gradient is taken through the
sampling distribution (\cref{prop:rb-reverse}). The two remaining pairings are forward KL on student
prefixes (for example GKD with forward KL on on-policy data) and reverse KL on teacher prefixes. These are
genuine discrepancies (\cref{prop:hybrid-div}) but do not equal any sequence-level $f$-divergence.
\end{remark}

\begin{proposition}[Hybrid argument for TV, and its converse]\label{prop:hybrid}
\Pv\Ck For every $x$,
\[
\TV(\Pseq,\Qseq)\le\sum_{t=1}^T\E_{s_t\sim\Pseq}\TV_t(s_t)
\quad\text{and}\quad
\TV(\Pseq,\Qseq)\le\sum_{t=1}^T\E_{s_t\sim\Qseq}\TV_t(s_t),
\]
where $\TV_t(s)=\TV(p_t(\cdot\mid s),q_t(\cdot\mid s))$. Conversely, for each $t$,
$\E_{s_t\sim\Pseq}\TV_t(s_t)\le2\,\TV(\Pseq,\Qseq)$, and the same holds with $\Qseq$.
\end{proposition}
\begin{proof}
Let $H_k$ draw $y_{1:k}$ from $p$ and $y_{k+1:T}$ from $q_\theta$, so that $H_T=\Pseq$ and $H_0=\Qseq$.
$H_k$ and $H_{k-1}$ share the law of $s_k$ and the continuation kernel after step $k$. By data processing,
$\TV(H_k,H_{k-1})\le\E_{s_k\sim\Pseq}\TV_k(s_k)$. Telescope. For the second bound, exchange the roles of
$p$ and $q_\theta$.
For the converse, let $\Pi^P_t$ and $\Pi^Q_t$ be the laws of $(s_t,y_t)$ and $\Pi^H:=\dteach t\otimes q_t$. Then
$\E_{\dteach t}\TV_t=\TV(\Pi^P_t,\Pi^H)\le\TV(\Pi^P_t,\Pi^Q_t)+\TV(\dstud t,\dteach t)\le2\,\TV(\Pseq,\Qseq)$,
because marginalisation cannot increase TV.
\end{proof}

\begin{proposition}[The mismatched pairings are divergences with a $\sqrt T$-lossy Pinsker bound]\label{prop:hybrid-div}
\Pv Let $\mathsf H_{\mathrm{on}}^{\KL}:=\sum_t\E_{s\sim\Qseq}\KL(p_t\|q_t)$ and
$\mathsf H_{\mathrm{off}}^{\mathrm{rKL}}:=\sum_t\E_{s\sim\Pseq}\KL(q_t\|p_t)$. Each is $\ge0$, and each
vanishes iff $\Pseq=\Qseq$. Moreover $\TV(\Pseq,\Qseq)\le\sqrt{T\,\mathsf H/2}$ for either one.
By contrast, the exact pairings give $\TV\le\sqrt{\KL/2}$ with $\KL$ the corresponding sequence-level KL.
The factor $T$ is an artefact of the proof, not a proven loss. A sequential Hellinger subadditivity
inequality \citep{foster2021statistical} gives $\TV\le C\sqrt{\mathsf H\log T}$ instead. Multi-start local search over binary
autoregressive pairs finds worst ratios $\TV^2/\mathsf H\approx0.50,0.60,0.69,0.74$ for $T=1,\dots,4$, far below $T/2$. These are
lower estimates of the true worst case.
\end{proposition}
\begin{proof}
If $\mathsf H_{\mathrm{on}}^{\KL}=0$, then $p_1=q_1$ at $x$. So $d_2^{\,p}=d_2^{\,\theta}$, hence $p_2=q_2$ on the
common support, and so on by induction. The bound follows from \cref{prop:hybrid}, Pinsker per token,
and Cauchy--Schwarz: $\sum_t\E\TV_t\le\sum_t\sqrt{e_t/2}\le\sqrt{T\sum_te_t/2}$.
\end{proof}

\begin{corollary}[Sequence-level TV from per-token errors]\label{cor:seq-tv}
\Pv
\[
\E_\cD\TV(P,Q_\theta)\le\min\Big\{\sqrt{T\eoff^{\KL}/2},\;\sqrt{1-e^{-T\eoff^{\KL}}},\;\sqrt{T\eon^{\mathrm{rKL}}/2},\;
T\eoff^{\TV},\;T\eon^{\TV}\Big\}.
\]
\end{corollary}
\begin{proof}
Combine \cref{thm:chain,prop:pinsker,prop:hybrid} with Jensen. The maps $k\mapsto\sqrt k$ and
$k\mapsto\sqrt{1-e^{-k}}$ are concave.
\end{proof}

\begin{example}[Neither route dominates]\label{ex:regimes}
\Pv
(a)~\emph{Diffuse errors.} Let $p_t=\mathrm{Bern}(\frac12)$ and $q_t=\mathrm{Bern}(\frac12+\delta)$ be i.i.d.
Then $T\eoff^{\TV}=T\delta$, whereas $\KL(P\|Q)=-\frac T2\log(1-4\delta^2)$, so the Pinsker route gives
$\TV\le\sqrt T\,\delta\,(1+o(1))$. This is the correct order, by the central limit theorem for the
log-likelihood ratio, and it improves on the hybrid route by $\sqrt T$.
(b)~\emph{Concentrated errors (deterministic teacher).} Let $p_t=\delta_{a}$ and $q_t(a)=1-\delta$. Then
$\TV=1-(1-\delta)^T\le T\delta$ and $\KL=-T\log(1-\delta)$, so $\sqrt{\KL/2}\approx\sqrt{T\delta/2}\gg T\delta$
when $T\delta\ll1$. Here the hybrid route is tight. In general, if $P$ is a point mass then
$\TV(P,Q)=1-e^{-\KL(P\|Q)}\le\KL(P\|Q)$: control is linear rather than square-root.
\end{example}

\begin{proposition}[Covariate shift]\label{prop:covshift}
\Pv For any measurable $g\ge0$ (for example a per-input divergence), under \cref{as:shift},
$\E_\mu g\le\rho\,\E_\cD g$. If $0\le g\le G$, then without any density-ratio assumption,
$|\E_\mu g-\E_\cD g|\le G\,\TV(\mu,\cD)$.
Bounded criteria ($\TV$, $\Hel^2$, $\JSD_\beta$) therefore transfer across input shift additively.
Unbounded ones (KL) need \cref{as:shift}.
\end{proposition}

\subsection{Which definition is the right one?}\label{sec:right}

\begin{definition}[Hierarchy of near-losslessness]\label{def:hierarchy}
At tolerance $\eps$ and deployment law $\mu$, the student is:
\begin{description}[nosep,font=\normalfont\itshape]
\item[behaviourally lossless] if $\E_\mu\TV(\Pseq,\Qseq)\le\eps$;
\item[likelihood lossless] if $\E_\mu\KL(\Pseq\|\Qseq)\le\eps$;
\item[two-sided lossless] if $\E_\mu[\KL(\Pseq\|\Qseq)+\KL(\Qseq\|\Pseq)]\le2\eps$;
\item[decoding lossless] if it is behaviourally lossless and the margin condition of
\cref{prop:margin} holds on the teacher's greedy paths (\cref{sec:decoding} quantifies this notion).
\end{description}
\end{definition}

\noindent\textbf{Position taken in this document.}
\begin{enumerate}[leftmargin=*]
\item \emph{The target is sequence-level, on the deployment law.} Behavioural losslessness,
$\E_\mu\TV(\Pseq,\Qseq)\le\eps$, is \emph{exactly} the statement that no bounded downstream task evaluated
by sampling changes by more than $\eps$ (\cref{prop:tv-task}). It is also exactly the statistical
indistinguishability of the two models from one sample. It is symmetric, bounded, and transfers across
covariate shift additively (\cref{prop:covshift}). Token-level criteria are \emph{surrogates}; they
relate to the target only through \cref{thm:chain,prop:hybrid}.
\item \emph{TV alone is not enough when probabilities are consumed or events are rare.} If the student's
likelihoods are used (scoring, calibration, uncertainty, perplexity), the relevant quantity is the
log-loss regret, which is \emph{exactly} $\E_\mu\KL(\Pseq\|\Qseq)$. If the student is sampled at scale and
rare events matter (hallucinations, unsafe outputs), the precision bound of \cref{prop:events}(ii--iii)
requires the \emph{reverse} direction. Hence two-sided losslessness. If decoding is greedy or beam
search, margins matter too (\cref{prop:margin}). Sampling-level losslessness then controls the decoded
output only through the teacher's probability of its own greedy prefix (\cref{thm:greedy-seq}), which
can be exponentially small.
\item \emph{Among trainable surrogates, sequence-level forward KL is canonical.} It dominates behavioural
loss via Pinsker and gives log-loss exactly. By \cref{thm:chain} its per-token decomposition uses
\emph{teacher prefixes only}, so it is estimable without student roll-outs and has no exposure-bias term
(\cref{sec:shift}). It is convex in the logits (\cref{prop:convexity}). More generally, KL is the only divergence
on the simplex that is both an $f$-divergence and a Bregman divergence \citep{amari2009alpha}. That
characterisation singles out the KL \emph{family}, not a direction: both argument orders qualify.
\item \emph{But when exactness is impossible, the surrogate must be chosen for the target.} Under
misspecification, minimisers of either KL direction can be arbitrarily bad in TV, whereas minimisers of
a bounded divergence such as $\JSD_\beta$ are TV-competitive (\cref{prop:agnostic}). This is why the
objective proposed in \cref{sec:objectives} is a \emph{sequence-level} $\JSD_\beta$ with an optional
forward-KL anchor.
\item \emph{In the truly near-lossless regime the question dissolves.} By \cref{prop:local}, all smooth
divergences agree to second order when likelihood ratios are bounded. The practically important choice
is then the \emph{prefix law} and the \emph{estimator}, not the divergence.
\end{enumerate}
In one sentence: a student should be called $\eps$-lossless if it is two-sided lossless at level $\eps$
under the deployment input law (behavioural losslessness at level $\sqrt{\eps/2}$ then follows). It
should be trained with a sequence-level divergence whose minimiser remains TV-competitive when
$\eps$-losslessness is out of reach.

